\documentclass[letterpaper,11pt]{article}

\usepackage{latexsym}
\usepackage[empty]{fullpage}
\usepackage{titlesec}
\usepackage{marvosym}
\usepackage[usenames,dvipsnames]{color}
\usepackage{verbatim}
\usepackage{enumitem}
\usepackage[pdftex]{hyperref}
\usepackage{fancyhdr}
\usepackage[utf8]{inputenc}
\usepackage[T2A]{fontenc}
\usepackage[english, russian]{babel}


\pagestyle{fancy}
\fancyhf{} % clear all header and footer fields
\fancyfoot{}
\renewcommand{\headrulewidth}{0pt}
\renewcommand{\footrulewidth}{0pt}

% Adjust margins
\addtolength{\oddsidemargin}{-0.375in}
\addtolength{\evensidemargin}{-0.375in}
\addtolength{\textwidth}{1in}
\addtolength{\topmargin}{-.5in}
\addtolength{\textheight}{1.0in}

\urlstyle{same}

\raggedbottom
\raggedright
\setlength{\tabcolsep}{0in}

% Sections formatting
\titleformat{\section}{
  \vspace{-4pt}\scshape\raggedright\large
}{}{0em}{}[\color{black}\titlerule \vspace{-5pt}]

%-------------------------
% Custom commands
\newcommand{\resumeItem}[2]{
  \item\small{
    \textbf{#1}{: #2 \vspace{-2pt}}
  }
}

\newcommand{\resumeSubheading}[4]{
  \vspace{-1pt}\item
    \begin{tabular*}{0.97\textwidth}{l@{\extracolsep{\fill}}r}
      \textbf{#1} & #2 \\
      \textit{\small#3} & \textit{\small #4} \\
    \end{tabular*}\vspace{-5pt}
}

\newcommand{\resumeSubItem}[2]{\resumeItem{#1}{#2}\vspace{-4pt}}

\renewcommand{\labelitemii}{$\circ$}

\newcommand{\resumeSubHeadingListStart}{\begin{itemize}[leftmargin=*]}
\newcommand{\resumeSubHeadingListEnd}{\end{itemize}}
\newcommand{\resumeItemListStart}{\begin{itemize}}
\newcommand{\resumeItemListEnd}{\end{itemize}\vspace{-5pt}}
\newcommand{\resumeProject}[1]{
  \vspace{-1pt}\item
    \textbf{#1}
    \vspace{-5pt}
}
%-------------------------------------------
%%%%%%  CV STARTS HERE  %%%%%%%%%%%%%%%%%%%%%%%%%%%%


\begin{document}

%----------HEADING-----------------
\noindent
\begin{tabular*}{\textwidth}{l@{\extracolsep{\fill}}r}
  \textbf{\Large Федор Фесенко} & Email: \href{mailto:fedorgym2005@gmail.com}{fedorgym2005@gmail.com} \\
  & tg: \href{https://t.me/howtozabotatshad}{@howtozabotatshad} \\
  & Github: \href{https://github.com/fedor13F}{fedor13F} \\
\end{tabular*}



%-----------Professional Summary-----------------
\section{Обо мне}
 \resumeItemListStart
    \item\small{Студент Сибирского федерального университета («Программная инженерия»), специализируюсь на ML с фокусом на RecSys и Information Retrieval. Основные интересы: RecSys, Information Retrieval, NLP.}
 \resumeItemListEnd


%-----------проекты-----------------
\section{Проекты}
\resumeSubHeadingListStart

\resumeProject{\href{https://github.com/fedor13F/Recommendation-System}{Рекомендательная система}}
\resumeItemListStart
  \resumeItem{Задача}{Построение end-to-end рекомендательной системы (кандидат-генерация + ранжирование) на данных Яндекс.Лавки в рамках соревнования от ШАД.}
  \resumeItem{Кандидат-генерация}{Реализовал и сравнил item-to-item, ALS и Two-Tower модель. Для Two-Tower исследовал две архитектуры: обучаемые эмбеддинги и Mean-энкодер (среднее эмбеддингов истории). Отбор кандидатов проводил по Recall@500, MRR@500 и Coverage.}
  \resumeItem{Обучение}{В кандгенах провёл эксперименты с BCE, Softmax, Sampled Softmax (uniform и in-batch негативы) и logQ-коррекцией для устранения bias. В ранкере сравнил лоссы YetiRank и PairLogit.}
  \resumeItem{Ранжирование}{Скоры кандидат-генераторов использовал как признаки для CatBoost. По Feature Importance, косинусы от ALS и Two-Tower вошли в топ-5. Итоговое качество: \textbf{NDCG@10 = 0.396}.}
\resumeItemListEnd

  \resumeProject{\href{https://github.com/fedor13F/Humanoid-BenchMark}{Humanoid Benchmark: метрика для оценки предвзятости LLM}}
\resumeItemListStart
  \resumeItem{Задача}{Разработка метрики для оценки изменения семантического bias LLM после дообучения на узком целевом домене, в рамках студкемпа \textbf{Яндекс Образования}.}
  \resumeItem{Метод}{Предложил метрику на основе \textbf{affect valence} — разности косинусных близостей эмбеддингов целевых слов к позитивным и негативным эмоциональным якорям. Для отделения локальных «фаворитов» от глобального сдвига домена ввёл \textbf{дисперсию valence} (\textbf{ΔVar}).}
  \resumeItem{Эксперименты}{Провёл domain-тюнинг Qwen2.5-1.5B и Qwen2.5-7B при разном объёме данных и числе эпох. Извлекал last hidden states трансформера, визуализировал через t-SNE.}
  \resumeItem{Результат}{Показал корреляцию метрики с объёмом дообучения (log(dataset\_size × epochs)).}
\resumeItemListEnd

\resumeSubHeadingListEnd


%-----------AWARDS-----------------
\section{Достижения}
 \resumeItemListStart
    \resumeItem{\href{https://contest.yandex.ru/contest/62615/standings/}{Призёр RuCode.START 2024}}{Чемпионат по алгоритмическому программированию (индивидуальный зачёт).}
 \resumeItemListEnd

%-----------EDUCATION-----------------
\section{Образование}
  \resumeSubHeadingListStart
    \resumeSubheading
      {Сибирский федеральный университет}{}
      {Программная инженерия, бакалавриат}{2023 -- н. в.}
  \resumeSubHeadingListEnd

%--------PROGRAMMING SKILLS------------
\section{Технологии и инструменты}
 \resumeSubHeadingListStart
   \item{\textbf{ML/DL}{: PyTorch, HuggingFace Transformers, scikit-learn, CatBoost, Optuna}}
   \item{\textbf{Data Tools}{: NumPy, Pandas, Matplotlib, Seaborn}}
   \item{\textbf{Languages}{: Python, SQL}}
   \item{\textbf{Infrastructure}{: Docker, FastAPI, MLflow, Airflow, Git}}
 \resumeSubHeadingListEnd



%-------------------------------------------
\end{document}